Un focus de llum és al punt $P(5,-1,5)$ i il·lumina perpendicularment una paret que ocupa part del pla
$\pi\colon 2x-y+2z-3=0$.

\begin{apartats}

\apartat[1,25]{0,75}
Escriu la recta perpendicular a la paret que passa per $P$ i troba en quin punt $Q$ la talla.

\begin{solucio}
Amb vector director $\vec n=(2,-1,2)$: $(5+2t,\ -1-t,\ 5+2t)$. Substituint:
$2(5+2t)-(-1-t)+2(5+2t)-3=18+9t=0$, $t=-2$, i $Q=(1,1,1)$.
\end{solucio}

\begin{tria}{focus}
\itemtria{simetric}{1}{1,25}
La paret és un mirall. Troba el punt $P'$ on es veu la imatge del focus, que és el simètric de $P$
respecte del pla.

\begin{solucio}
$Q$ és el punt mitjà de $P$ i $P'$: $P'=2Q-P=(2-5,\ 2+1,\ 2-5)=(-3,3,-3)$.
\end{solucio}

\itemtria{distancia-dues-maneres}{1}{1,25}
Calcula la distància del focus a la paret amb la fórmula de la distància d'un punt a un pla, i també
com a distància entre $P$ i $Q$.

\begin{solucio}
$d(P,\pi)=\dfrac{|10+1+10-3|}{\sqrt{4+1+4}}=\dfrac{18}3=6$. D'altra banda,
$\overrightarrow{PQ}=(-4,2,-4)$ i $|\overrightarrow{PQ}|=\sqrt{16+4+16}=6$.
\end{solucio}
\end{tria}

\begin{nomesllarg}
\apartat{0,75}
Troba el pla paral·lel a la paret que conté el focus, i la distància entre aquest pla i la paret.

\begin{solucio}
$2x-y+2z+D=0$ amb $10+1+10+D=0$, $D=-21$: el pla $2x-y+2z-21=0$. La distància entre els dos plans, amb el
punt $Q(1,1,1)$ de la paret, és $\dfrac{|2-1+2-21|}{3}=\dfrac{18}{3}=6$.
\end{solucio}
\end{nomesllarg}

\end{apartats}
